% TOP_PROBLEM: {"id": 30007088, "title": "Shareshian–Wachs elementary positivity for chromatic quasisymmetric functions", "category_id": 2, "category": {"id": 2, "name": "combinatorics", "display_name": "Combinatorics", "description": "Counting problems, graph theory, discrete structures.", "slug": "combinatorics", "order_index": 2, "created_at": "2026-07-31T15:26:25.671Z"}, "status": "open"}
% FIRST_POSTED: 2026-10-06
% ATTEMPT_STATUS: solved
% OPENAI_CLAIM: {"schema_version":1,"source_repo":"https://github.com/openai/math","source_commit":"adc7f1241b42e322a6451854ab7e4b4c146bf78a","release_date":"2026-10-06","match":"full","resolution":"proved","independently_reviewed":false,"second_pass":"agreed","families":[{"family":"169","title":"Shareshian–Wachs elementary positivity","subject":"Combinatorics","manuscripts":[{"title":"Elementary positivity of chromatic quasisymmetric functions","pdf_path":"preprints/Elementary-Positivity-of-Chromatic-Quasisymmetric-Functions-September-24-2026/paper.pdf","readme_path":"preprints/Elementary-Positivity-of-Chromatic-Quasisymmetric-Functions-September-24-2026/README.md","date":"2026-09-24","theorem_ref":"preprints/Elementary-Positivity-of-Chromatic-Quasisymmetric-Functions-September-24-2026/build/sections/01-introduction.tex#main:theorem"}],"lean":{"doc_path":"lean/docs/169.md","comparators":["lean/ComparatorChallenges/ElementaryPositivity.json","lean/ComparatorChallenges/ElementaryPositivity.lean"],"declarations":[],"covers_main_theorem":true},"reasoning_trace":null}]}
\documentclass[11pt]{article}
\usepackage[margin=25mm]{geometry}
\usepackage{fontspec}
\setmainfont{Latin Modern Roman}
\usepackage{amsmath,amssymb,mathtools}
\usepackage{unicode-math}
\setmathfont{Latin Modern Math}
\usepackage{enumitem}
\usepackage{xurl,hyperref,bookmark}
\hypersetup{colorlinks=true,urlcolor=blue}
\setcounter{secnumdepth}{0}
\setlength{\parindent}{0pt}
\setlength{\parskip}{5pt}
\setlength{\emergencystretch}{3em}
\begin{document}
\section*{Shareshian–Wachs elementary positivity for chromatic quasisymmetric functions}

\subsection{OpenAI's claim}
\textit{Elementary positivity for unit interval graphs}: OpenAI claims nonnegative $\mathbb N[q]$ coefficients in the elementary-basis expansion of chromatic quasisymmetric functions of natural unit interval graphs. This concerns the elementary-positivity portion of the Shareshian--Wachs conjecture.

\subsection{Outline of the argument}
\textit{Elementary positivity for unit interval graphs}: Positive wall transformations produce sections indexed by exponent-lattice elements. Interval rotations identify the relevant graph functions with single positive sections, and a support argument forces the contributing exponents onto distinguished anchors. A packet bijection then extracts a nonnegative combinatorial expansion whose powers of $q$ record the accumulated jumps.

\subsection{Scope relative to this problem}
\textit{Elementary positivity for unit interval graphs}: Theorem main:theorem gives a formal symmetric-function identity with each elementary coefficient in N[q] for every natural unit interval graph on its natural labels. It proves the graded statement, rather than only its q=1 specialization, and includes an explicit permutation witness.

With the hypotheses and conventions above, the claimed result addresses the whole mathematical assertion on this page.

\subsection{Formal verification}
Family 169 describes Lean coverage for a nonnegative elementary expansion for natural unit interval graphs with its inversion weights.

\textit{Elementary positivity for unit interval graphs}: The manuscript links \texttt{ElementaryPositivity.lean}.

The documented selection includes this manuscript’s main assertion; that assertion retains the restrictions stated above.

This site has not run the Lean development. These links describe OpenAI’s selected formal statements and should be read with their stated scope.

\subsection{Status caveat}
OpenAI reports that an internal model produced the manuscripts. Their arguments have not been independently verified as proofs by this site. Any portion outside the documented formal scope may contain mistakes, and the catalogue label reports the announced mathematical claim.

\subsection{Sources}
\begin{itemize}
\item OpenAI, \emph{Elementary positivity of chromatic quasisymmetric functions}, 2026-09-24; citation key \texttt{OAI:Elementary-Positivity-of-Chromatic-Quasisymmetric-Functions-September-24-2026}. \href{https://github.com/openai/math/blob/adc7f1241b42e322a6451854ab7e4b4c146bf78a/preprints/Elementary-Positivity-of-Chromatic-Quasisymmetric-Functions-September-24-2026/paper.pdf}{PDF (GitHub)}; \href{https://github.com/openai/math/raw/adc7f1241b42e322a6451854ab7e4b4c146bf78a/preprints/Elementary-Positivity-of-Chromatic-Quasisymmetric-Functions-September-24-2026/paper.pdf}{PDF (direct)}; \href{https://github.com/openai/math/blob/adc7f1241b42e322a6451854ab7e4b4c146bf78a/preprints/Elementary-Positivity-of-Chromatic-Quasisymmetric-Functions-September-24-2026/README.md}{Manuscript page and citation}; \href{https://github.com/openai/math/blob/main/preprints/Elementary-Positivity-of-Chromatic-Quasisymmetric-Functions-September-24-2026/paper.pdf}{Latest version}.
\item Family 169: \href{https://github.com/openai/math/blob/adc7f1241b42e322a6451854ab7e4b4c146bf78a/lean/docs/169.md}{Lean scope}.
\item \href{https://github.com/openai/math/blob/adc7f1241b42e322a6451854ab7e4b4c146bf78a/lean/ComparatorChallenges/ElementaryPositivity.json}{Comparator ElementaryPositivity.json}.
\item \href{https://github.com/openai/math/blob/adc7f1241b42e322a6451854ab7e4b4c146bf78a/lean/ComparatorChallenges/ElementaryPositivity.lean}{Comparator ElementaryPositivity.lean}.
\item \href{https://github.com/openai/math/blob/adc7f1241b42e322a6451854ab7e4b4c146bf78a/CONTENTS.md}{OpenAI catalogue entry}.
\end{itemize}
\end{document}
